\chapter{Fundamental Theorem of Algebra}

\section{Polynomial roots as computable objects}

\begin{definition}[Rational-complex polynomial]\label{def:qcomplex-poly}
\lean{ComputableAnalysis.CPoly.Coeffs}
\uses{def:complex-raw}
A rational-complex polynomial is a finite list
\[
  (a_0,\ldots,a_d),\qquad a_j\in\Qi,
\]
interpreted as
\[
  p(z)=a_0+a_1z+\cdots+a_dz^d.
\]
Positive degree means that some coefficient with index \(d\ge 1\) is
nonzero.
\end{definition}

\begin{definition}[Computable root]\label{def:computable-root}
\lean{ComputableAnalysis.IsComputableRoot}
\uses{def:qcomplex-poly,def:complex-raw,def:raw-equiv}
Let \(p\in\Qi[z]\).  A computable root of \(p\) is a complex number \(z\)
such that the interval evaluation of \(p(z)\) is equivalent to zero.
\end{definition}

\begin{definition}[Algebraic complex number]\label{def:algebraic-complex}
\lean{ComputableAnalysis.AlgebraicComplex}
\uses{def:qcomplex-poly,def:computable-root}
An algebraic complex number is a computable complex number together with a
rational-complex annihilating polynomial.  The intended mature version also
stores enough isolation data to make arithmetic and equality computable.
\end{definition}

\section{Base cases}

\begin{theorem}[Monic linear polynomials]\label{thm:monic-linear-root}
\lean{ComputableAnalysis.monicLinear_has_computable_root}
\uses{def:computable-root}
\leanok
For every \(r\in\Qi\), the polynomial \(z-r\) has the computable root \(r\).
\end{theorem}

\begin{theorem}[The root of \(z^2+1\)]\label{thm:z2-plus-one-root}
\lean{ComputableAnalysis.zSqPlusOne_has_computable_root}
\uses{def:computable-root}
\leanok
The polynomial \(z^2+1\) has \(i\) as a computable root:
\[
  i^2+1=0.
\]
\end{theorem}

\section{FTA statements}

\begin{conjecture}[Computable FTA]\label{conj:computable-fta}
\lean{ComputableAnalysis.ComputableFTA}
\uses{def:qcomplex-poly,def:computable-root}
Every positive-degree rational-complex polynomial has a computable complex
root.  In Lean this means that the root candidate is a valid complex raw real
and the interval evaluation of the polynomial at that candidate represents
zero.
\end{conjecture}

\begin{conjecture}[Algebraic FTA]\label{conj:algebraic-fta}
\lean{ComputableAnalysis.AlgebraicFTA}
\uses{def:qcomplex-poly,def:algebraic-complex}
Every positive-degree rational-complex polynomial has an algebraic complex
root whose polynomial evaluation represents zero.
\end{conjecture}

\begin{theorem}[Algebraic roots give computable roots]
\label{thm:computable-fta-of-algebraic-fta}
\lean{ComputableAnalysis.ComputableFTA_of_AlgebraicFTA}
\uses{conj:computable-fta,conj:algebraic-fta}
\leanok
The algebraic FTA implies the computable FTA.
\end{theorem}

\section{Two constructive routes}

\begin{proposition}[Algebraic route]\label{prop:fta-algebraic-route}
\uses{conj:algebraic-fta}
It is enough to build an algebraic closure of \(\Qi\) with computable
embeddings into complex interval boxes.  The finite proof obligations are:
\begin{enumerate}
  \item factor a nonconstant polynomial into irreducible factors over \(\Qi\);
  \item adjoin one formal root of an irreducible factor by a quotient algebra;
  \item isolate a complex embedding of that root by rational boxes;
  \item prove that interval evaluation overlaps zero.
\end{enumerate}
\end{proposition}

\begin{proposition}[Analytic search route]\label{prop:fta-analytic-route}
\uses{conj:computable-fta,def:qbox}
An alternative proof avoids abstract algebraic closure.  For a positive-degree
polynomial \(p\), choose a rational square \(B\subset\C\) so large that
\[
  |a_d z^d| >
  |a_{d-1}z^{d-1}+\cdots+a_0|
  \qquad (z\in\partial B).
\]
Then the winding number of \(p(\partial B)\) around zero is \(d\).  Subdivide
\(B\) rationally, using interval arithmetic to discard boxes whose image
does not meet zero and to keep a box whose root count is positive.  A nested
sequence of such boxes is a computable root.
\end{proposition}

\begin{remark}
The algebraic route is closer to pre-completeness algebra.  The analytic
route is closer to the usual complex-analytic proof of FTA, but it must be
made finite by replacing open-cover and compactness arguments with explicit
box bounds and winding-number computations.
\end{remark}
